% ==============================================================================
% Plantilla Institucional de Trabajo Terminal — UPIIZ IPN
% Archivo: capitulos/02-justificacion.tex
% ==============================================================================
% LÍMITE INSTITUCIONAL: Máximo 1 página.
% ==============================================================================

\chapter{Justificación}
\label{cap:justificacion}

La investigación y desarrollo en sistemas vehiculares autónomos y asistidos demanda entornos de experimentación seguros, accesibles y altamente reproducibles. La experimentación directa sobre vehículos de tamaño real conlleva elevados costos de instrumentación, requerimientos de infraestructura especializada y riesgos considerables para la integridad física de los operadores e investigadores. En este contexto, el uso de plataformas mecatrónicas a escala 1:10 constituye una alternativa técnica idónea para la validación rigurosa de algoritmos de control, estimación de estados y procesamiento digital en tiempo real.

Desde la perspectiva técnica, el control de la velocidad longitudinal y del ángulo de dirección involucra dinámicas electromecánicas acopladas, no linealidades inherentes a la tracción y efectos de fricción e inercia variables. La formulación del sistema en espacio de estados proporciona un marco analítico robusto que permite no solo sintetizar leyes de control multivariable con especificaciones dinámicas estrictas (tiempo de establecimiento, sobrepaso y error en régimen permanente), sino también evaluar formalmente la controlabilidad y observabilidad del sistema y estimar variables no medibles directamente mediante observadores o filtros digitales.

Institucionalmente, el presente proyecto se encuentra plenamente alineado con la \textbf{Línea de investigación IV: Diseño e implementación de sistemas o técnicas de control} del Programa Académico de Ingeniería Mecatrónica de la UPIIZ-IPN. Si bien el énfasis principal radica en la teoría y aplicación de control moderno, el proyecto exige la integración sinérgica de las cuatro áreas de la disciplina: el análisis cinemático y geométrico del chasis (Mecánica), el diseño de etapas de potencia y acondicionamiento sensorial (Electrónica), el modelado y síntesis de la ley de control (Control) y el desarrollo de firmware embebido determinista (Programación).

Socialmente, el proyecto aporta una metodología transferible para el diseño de bancos de prueba docentes y de investigación en el Instituto Politécnico Nacional, facilitando la formación práctica de ingenieros en el desarrollo de tecnologías de movilidad terrestre y sistemas de transporte inteligente.
